% !PRINTED 659
\weberpage{659}

%<<CONT
\weq{1}{\eta = \frac{\mu}{\nu},}

and $\mu$ is called the numerator, $\nu$ the denominator of $\eta$. These
two functions are not completely determined by $\eta$; what is determined,
however, are the functionals that remain when all common functional factors
are removed from numerator and denominator. For if

\weq{2}{\frac{\mu}{\nu} = \frac{\mu'}{\nu'}; \quad \mu\nu' = \nu\mu',}

then every prime functional that divides $\nu$ but not $\mu$ must divide
$\nu'$. The functionals $\bar\alpha$ and $\bar\beta$ thus determined from
$\eta$ we call the numerator functional and denominator functional of $\eta$.
In order to represent $\eta$ by functions, one first takes a function $\nu$
divisible by $\bar\beta$ but otherwise arbitrary, and sets
$\nu = \bar\beta\bar\gamma$; in this $\bar\gamma$ may be assumed relatively
prime to any functional one pleases. Then $\eta\nu$ is an integral function
divisible by $\bar\gamma$, and one sets

\[ \eta\nu = \mu = \bar\alpha\bar\gamma. \]

The functional $\bar\alpha$ cannot, however, as we shall prove later on, be
chosen arbitrarily any longer once $\bar\beta$ is given.

The functionals $\bar\alpha, \bar\beta$ we shall also call, for short, \gesperrt{the
numerator and the denominator} of $\eta$.

If $\pi$ is a prime functional, then there are \gesperrt{numbers} (constants) $a, b,
a', b'$ which satisfy the congruences

\[ \mu \equiv a,\ \ \nu \equiv b,\ \ \mu' \equiv a',\ \ \nu' \equiv b' \quad (\mathrm{mod}\ \pi) \]

and from (2) it follows that

\[ ab' = ba', \quad \frac{a}{b} = \frac{a'}{b'} = c, \]

provided that $\nu$ and $\nu'$ are not divisible by $\pi$. Then $\eta - c$ is
a function in $\Omega$ in which the numerator, but not the denominator, is
divisible by $\pi$. If $\varrho$ is an integral or fractional function in
which the numerator is divisible by $\pi$ but not by $\pi^2$, and the
denominator is likewise not divisible by $\pi$, then $(\eta - c) : \varrho$
is a function whose denominator is not divisible by $\pi$ and whose numerator
contains the factor $\pi$ one time fewer than $(\eta - c)$. We set:

\[ \eta - c = \varrho\eta_1. \]

We apply the same consideration to $\eta_1$, and so obtain a sequence of
equations that may be continued as far as we please:
%CONT>>
